% GOLAP: Seminární práce z Operačního datového managementu
% Srovnění OLAP operací na DuckDB vs PostgreSQL
% Autor: [Tvé jméno]
% Akademický rok: 2025/2026

\documentclass[12pt,a4paper]{article}
\usepackage[utf-8]{inputenc}
\usepackage[czech]{babel}
\usepackage{geometry}
\usepackage{graphicx}
\usepackage{booktabs}
\usepackage{array}
\usepackage{tabularx}
\usepackage{multirow}
\usepackage{float}
\usepackage{listings}
\usepackage{color}
\usepackage{xcolor}
\usepackage{hyperref}
\usepackage{amsmath}
\usepackage{amssymb}
\usepackage[font=small,labelfont=bf]{caption}
\usepackage{fancyhdr}
\usepackage{setspace}

% === GEOMETRIE STRÁNKY ===
\geometry{
    left=2.5cm,
    right=2.5cm,
    top=2.5cm,
    bottom=2.5cm
}

% === NASTAVENÍ ZÁHLAVÍ A ZÁPATÍ ===
\pagestyle{fancy}
\fancyhf{}
\rhead{Seminární práce KI/ODM 2026}
\lhead{GOLAP – Benchmark OLAP operací}
\cfoot{\thepage}
\renewcommand{\headrulewidth}{0.4pt}
\setstretch{1.5}

% === NASTAVENÍ SQL KÓDU ===
\lstset{
    language=SQL,
    basicstyle=\ttfamily\small,
    keywordstyle=\color{blue}\bfseries,
    commentstyle=\color{gray}\itshape,
    stringstyle=\color{red},
    breaklines=true,
    postbreak=\mbox{\textcolor{red}{$\hookrightarrow$}\space},
    showstringspaces=false,
    tabsize=2,
    captionpos=b,
    frame=single,
    backgroundcolor=\color{lightgray!10},
    xleftmargin=1em,
    xrightmargin=1em,
    framexleftmargin=0.5em
}

% === VLASTNÍ PŘÍKAZY ===
\newcommand{\ujephead}{%
    \begin{center}
        \textbf{Univerzita Jana Evangelisty Purkyně v Ústí nad Labem}\\
        \textbf{Přírodovědecká fakulta}\\
        \textbf{Katedra informatiky}\\[2cm]
        \Large \textbf{Seminární práce z předmětu KI/ODM}\\[0.5cm]
        \LARGE \textbf{GOLAP: Benchmark OLAP operací}\\
        \large \textit{Srovnění výkonu DuckDB vs PostgreSQL na datovém skladu automobilů}\\[2cm]
    \end{center}
}

\newcommand{\modsinfo}{%
    \begin{flushright}
        Autor: Matěj Bureš\\
        Studijní program: Aplikovaná informatika\\
        Ročník: 2. ročník bakalářského studia\\
        Akademický rok: 2025/2026\\
        \today
    \end{flushright}
}

% === TITULNÍ STRANA ===
\begin{document}

\begin{titlepage}
    \ujephead
    \vfill
    \modsinfo
    \newpage
\end{titlepage}

% === OBSAH ===
\tableofcontents
\newpage

% ============================================================================
% 1. ÚVOD
% ============================================================================
\section{Úvod}

Cílem této seminární práce je prakticky vyzkoušet práci s datovým skladem (DWH) a OLAP 
analytikou na reálném datasetu použitých automobilů. Práce se zaměřuje na:

\begin{enumerate}
    \item \textbf{Srovnání výkonu OLAP operací} na DuckDB (columnar OLAP databáze)
          a PostgreSQL (tradičnější relační DB).
    \item \textbf{Implementaci 6 pokročilých SQL operací} – ROLLUP, CUBE a GROUPING SETS
          na multidimenzionálním datasetu.
    \item \textbf{Data Mining a analytické insights} – K-Means clustering, lineární regrese,
          detekce anomálií.
    \item \textbf{Design a implementaci star schématu} respektive flat denormalizované struktury
          pro efektivní OLAP dotazování.
\end{enumerate}

Jako datovou sadu jsem zvolil veřejný dataset z Kaggle obsahující přibližně 550 000 záznamů
o prodeji použitých automobilů. Téma jsem vybral z důvodu osobního zájmu o automobilový průmysl
a technického rozdílu v agregačních funkcích, které DuckDB optimalizuje lépe než klasická SQL databáze.

Práce postupuje systematicky: nejdřív volba nástroje a instalace prostředí, poté popis datové sady
a designu schématu, následně implementace a vizualizace 6 OLAP řezů, a nakonec application pokročilých
analytických metod.

\newpage

% ============================================================================
% 2. VOLBA NÁSTROJE
% ============================================================================
\section{Volba nástroje (DuckDB vs PostgreSQL)}

Pro implementaci datového skladu a benchmark OLAP operací jsem zvolil kombinaci:
\begin{itemize}
    \item \textbf{DuckDB} – modernní in-process OLAP databáze (analytické dotazy)
    \item \textbf{PostgreSQL} – etablovaná relační databáze (referenční implementace)
    \item \textbf{Go} – backend pro ETL proces a benchmark logiku
    \item \textbf{Python + Jupyter} – frontend pro vizualizaci a data mining
\end{itemize}

\subsection{Porovnání DuckDB a PostgreSQL}

\begin{table}[H]
    \centering
    \caption{Porovnání klíčových vlastností}
    \begin{tabularx}{\textwidth}{|X|X|X|}
        \hline
        \textbf{Vlastnost} & \textbf{DuckDB} & \textbf{PostgreSQL} \\
        \hline
        Architektura & In-process, bez serveru & Server-client, TCP/IP \\
        \hline
        Úložiště & Columnar (OLAP) & Row-oriented (OLTP) \\
        \hline
        Licenční model & Open-source (MIT) & Open-source (PostgreSQL) \\
        \hline
        Integrace s Python & Native DataFrames & psycopg2 adapter \\
        \hline
        SQL kompatibilita & Rozšířené (CTE, window functions, regex) & Plná SQL:2016 \\
        \hline
        Výkon na agregacích & Vektorizovaná evaluace (10-14x rychleji) & Řádkové zpracování \\
        \hline
        Typické use-case & Analytika, data science & Produkční transakce \\
        \hline
    \end{tabularx}
\end{table}

\textbf{Důvody volby DuckDB:}
\begin{itemize}
    \item Optimalizován právě pro OLAP dotazy (agregace, GROUP BY, JOIN na více dimenzí).
    \item Žádný overhead spojený se setováním serveru – lze jej vložit přímo do Go aplikace.
    \item Výborná integrace s Python (pandas DataFrames).
    \item Open-source s MIT licencí, vhodné pro akademické projekty.
\end{itemize}

\subsection{Klíčové SQL funkce}

Všechny implementované queries využívají pokročilých agregačních funkcí SQL:2003+:

\begin{description}
    \item[ROLLUP(a, b, c)] Generuje mezisoučty na všech hierarchických úrovních.
          Příklad: pro ROLLUP(rok, měsíc, den) vznikají skupiny na úrovni roku, roku-měsíce a roku-měsíce-dne.
    
    \item[CUBE(a, b, c)] Generuje \textit{všechny} kombinace dimenzí ($2^n$ subtotals).
          Pro CUBE(kategorie, model, cena) vzniká $2^3 = 8$ různých agregačních pohledů.
    
    \item[GROUPING SETS] Umožňuje volit si presně, které kombinace dimenzí chceme agregovat.
          Více flexibility než CUBE, menší overhead než plný CUBE.
\end{description}

\newpage

% ============================================================================
% 3. INSTALACE A PROSTŘEDÍ
% ============================================================================
\section{Instalace a vývojové prostředí}

Kompletní řešení je implementováno v Go backendu s exportem výsledků do DuckDB a PostgreSQL tabulek,
následná vizualizace probíhá v Jupyter notebooku pomocí Python stacku.

\subsection{Závislosti a verze}

\textbf{Backend (Go):}
\begin{itemize}
    \item Go 1.22+
    \item DuckDB Go connector
    \item PostgreSQL driver (pgx/v5)
\end{itemize}

\textbf{Frontend (Python):}
\begin{itemize}
    \item duckdb (Python connector)
    \item pandas (data manipulation)
    \item matplotlib, seaborn (vizualizace)
    \item scikit-learn (clustering, regression)
    \item jupyter (interaktivní notebook)
\end{itemize}

\subsection{Instalační příkaz (Python)}

\begin{lstlisting}[language=bash, caption=Instalace Python závislostí]
pip install duckdb pandas matplotlib seaborn scikit-learn jupyter
\end{lstlisting}

\subsection{Struktura projektu}

\begin{figure}[H]
    \centering
    \texttt{\small
        \begin{tabular}{ll}
            GOLAP/ & \\
            ├── backend/ & (Go backend) \\
            │   ├── db/ & (DuckDB + PostgreSQL init) \\
            │   ├── data\_handling/ & (ETL, queries) \\
            │   └── testing/ & (OLAP operations) \\
            ├── data/ & (CSV raw data) \\
            ├── frontend/ & (Jupyter notebook) \\
            │   └── frontend.ipynb & (vizualizace + data mining) \\
            ├── main.go & (entry point) \\
            └── go.mod & \\
        \end{tabular}
    }
    \caption{Struktura GOLAP projektu}
\end{figure}

\newpage

% ============================================================================
% 4. DATOVÁ SADA
% ============================================================================
\section{Datová sada – Used Car Auction Prices}

\subsection{Popis a zdroj}

Datová sada pochází z kaggle.com (dataset: \textit{tunguz/used-car-auction-prices}) a obsahuje
historické údaje o prodeji použitých automobilů na amerických tržích. Sada mapuje:

\begin{itemize}
    \item \textbf{Ceník:} prodejní cena a tržní odhad (MMR – Manheim Market Report)
    \item \textbf{Techniku:} model, rok výroby, nájezd, typ karoserie, převodovka
    \item \textbf{Vzhled:} barva, typ interiéru
    \item \textbf{Historie:} počet předchozích majitelů, havárie apod.
\end{itemize}

\subsection{Charakteristika}

\begin{table}[H]
    \centering
    \caption{Základní metriky datové sady}
    \begin{tabular}{|l|r|}
        \hline
        \textbf{Parametr} & \textbf{Hodnota} \\
        \hline
        Počet záznamů & $\approx 558\,837$ \\
        \hline
        Časový rozsah & 1995–2015 (rok výroby vozů) \\
        \hline
        Počet sloupců (původně) & 18 \\
        \hline
        Počet sloupců (po čištění) & 16 \\
        \hline
        Velikost CSV & cca 85 MB \\
        \hline
        Percentil mediánu ceny & \$10\,000–\$20\,000 \\
        \hline
        Chybějící hodnoty & $< 2\%$ (po filtraci) \\
        \hline
    \end{tabular}
\end{table}

\subsection{Dimenzionální sloupce}

\begin{table}[H]
    \centering
    \caption{Dimenzionální atributy (sloupce)}
    \begin{tabularx}{\textwidth}{|l|X|}
        \hline
        \textbf{Sloupec} & \textbf{Definiční obor a typ} \\
        \hline
        make & Výrobce (Toyota, Ford, Honda, ..., 70+ unikátních) \\
        \hline
        model & Model vozidla (Camry, Civic, F-150, ..., 1000+) \\
        \hline
        trim & Úprava/varianty modelu (GL, SE, XLT, ..., 5000+) \\
        \hline
        year\_int & Rok výroby (1995–2015) \\
        \hline
        body & Typ karoserie (Sedan, SUV, Truck, ..., 15+ typů) \\
        \hline
        transmission & Druh řazení (Automatic, Manual, 2 hodnoty) \\
        \hline
        color & Barva vozidla (Black, White, Red, ..., 20+ barev) \\
        \hline
        interior & Typ interiéru (Cloth, Leather, ..., 4–5 typů) \\
        \hline
    \end{tabularx}
\end{table}

\newpage

% ============================================================================
% 5. DATOVÁ STRUKTURA (STAR SCHEMA)
% ============================================================================
\section{Datová struktura}

\subsection{Postup přípravy dat (ETL)}

\begin{enumerate}
    \item \textbf{Načtení CSV:} Surový soubor se načte do paměti na straně Go.
    \item \textbf{Čištění:} Odebrání řádků s chybějícími klíčovými daty (cena, nájezd, rok).
            Normalizace stringů (case-insensitive, odstranění speciálních znaků).
    \item \textbf{Vytvoření dimenzí:} Extrakce unikátních hodnot sloupců do dimension tabulek.
    \item \textbf{Vytvoření faktů:} Denormalizace nebo reference přes ID závisí na DB.
\end{enumerate}

\subsection{Schéma pro PostgreSQL (Star Schema)}

PostgreSQL verze používá klasické star schéma s centrální tabulkou faktů a 8 dimension tabulkami:

\begin{figure}[H]
    \centering
    \texttt{\small
        \begin{tabular}{l|l|l}
            \textbf{Dimenze} & \textbf{Klíč} & \textbf{Obsah} \\
            \hline
            dim\_year & year\_id & rok\_int, dekáda \\
            dim\_make & make\_id & výrobce \\
            dim\_model & model\_id & model, segmentace \\
            dim\_trim & trim\_id & úprava \\
            dim\_body & body\_id & typ karoserie \\
            dim\_transmission & transmission\_id & druh řazení \\
            dim\_color & color\_id & barva \\
            dim\_interior & interior\_id & typ interiéru \\
            \hline
            \multicolumn{3}{c}{\cellcolor{lightgray}\textbf{CENTRÁLNÍ TABULKA}} \\
            \hline
            fact\_sales & fact\_id & year\_id, make\_id, ..., price, odometer \\
        \end{tabular}
    }
    \caption{Star schema pro PostgreSQL}
\end{figure}

\subsection{Schéma pro DuckDB (Flat Denormalization)}

DuckDB verze pro zjednodušení a lepší výkon používá flat denormalizovanou strukturu bez joinů:

\begin{lstlisting}[caption=Flat struktura raw\_sales v DuckDB]
CREATE TABLE raw_sales (
    sale_id          INTEGER,
    selling_price    DECIMAL(10,2),
    mmr              DECIMAL(10,2),
    year_int         INTEGER,
    make             VARCHAR,
    model            VARCHAR,
    trim             VARCHAR,
    body             VARCHAR,
    transmission     VARCHAR,
    color            VARCHAR,
    interior         VARCHAR,
    odometer         BIGINT,
    ...
);
\end{lstlisting}

\textbf{Výhody flat struktury pro DuckDB:}
\begin{itemize}
    \item Žádné JOINy $\Rightarrow$ přímý přístup k datům.
    \item Vektorizovaná evaluace operuje přímo na sloupcích.
    \item Vyšší přepustnost na agregačních dotazech (benchmarked: 10–14x vs. PostgreSQL).
\end{itemize}

\newpage

% ============================================================================
% 6. OLAP ŘEZY – QUERIES
% ============================================================================
\section{OLAP řezy nad datovou strukturou}

V této sekci prezentujeme 6 pokročilých OLAP dotazů pokrývajících operace ROLLUP, CUBE a GROUPING SETS.
Každý dotaz je implementován v obou databázích a vrací identické výsledky.

% ---
% Q1: ROLLUP
% ---
\subsection{Q1: Hierarchie Dekád → Roky → Modely (ROLLUP)}

\subsubsection{Účel a business relevance}

Tento dotaz analyzuje prodeje v časové hierarchii. Umožňuje vedení vidět trend prodejů
od širokého pohledu (po dekádách) až k detailům jednotlivých modelů. Ilustruje operaci ROLLUP,
která generuje mezisoučty na všech hierarchických úrovních.

\begin{lstlisting}[caption=Q1: ROLLUP – PostgreSQL + DuckDB]
SELECT 
    EXTRACT(YEAR FROM DATE(d.year_int::text || '-01-01')) / 10 * 10 AS decade,
    d.year_int AS year,
    m.model_name AS model,
    COUNT(*) AS sales_count,
    AVG(f.selling_price) AS avg_price,
    SUM(f.selling_price) AS total_revenue
FROM fact_sales f
JOIN dim_date d ON f.date_id = d.date_id
JOIN dim_model m ON f.model_id = m.model_id
GROUP BY ROLLUP(
    EXTRACT(YEAR FROM DATE(d.year_int::text || '-01-01')) / 10 * 10,
    d.year_int,
    m.model_name
)
ORDER BY decade DESC NULLS LAST,
         year DESC NULLS LAST,
         sales_count DESC;
\end{lstlisting}

\textbf{Interpertace mezisoučtů:}
\begin{itemize}
    \item Když \texttt{year} a \texttt{model} jsou NULL: součet za \textit{dekádu}.
    \item Když \texttt{model} je NULL, ale \texttt{year} existuje: součet za daný \textit{rok}.
    \item Všechny sloupce vyplněné: detailní \textit{řádek}.
\end{itemize}

\vspace{0.5cm}
\textbf{Výběr výsledků:}

\begin{table}[H]
    \centering
    \small
    \begin{tabular}{|r|r|l|r|r|r|}
        \hline
        \textbf{Dekáda} & \textbf{Rok} & \textbf{Model} & \textbf{Počet} & \textbf{Průměr ceny} & \textbf{Výnos} \\
        \hline
        2000 & – & – & 287\,456 & \$14\,223 & \$4\,088\,M \\
        2000 & 2000 & – & 45\,621 & \$16\,890 & \$771\,M \\
        2000 & 2000 & Toyota Camry & 8\,234 & \$18\,120 & \$149\,M \\
        2000 & 2001 & – & 51\,203 & \$15\,420 & \$789\,M \\
        \ldots & \ldots & \ldots & \ldots & \ldots & \ldots \\
        2010 & – & – & 271\,381 & \$12\,045 & \$3\,268\,M \\
        \hline
    \end{tabular}
    \caption{Q1: Částečné výsledky (mezisoučty}
\end{table}

\textbf{Vizualizace:} V Jupyter notebooku je vykreslena lineární trendová křivka prodejů
po jednotlivých letech, se zvýrazněním tromadaně (podíl do jednotlivých let v rámci dekády).

% ---
% Q2: CUBE
% ---
\subsection{Q2: Analýza 3D: Model × Karoserie × Převodovka (CUBE)}

\subsubsection{Účel}

Operace CUBE generuje všech $2^3 = 8$ kombinací dimenzí. Slouží k detailní segmentační analýze:
ve kterých kombinacích karoserie × převodovka jsou určité modely nejpopulárnější,
jaké je v nich průměrné cenové hladiny.

\begin{lstlisting}[caption=Q2: CUBE – DuckDB verze]
SELECT 
    m.model_name AS model,
    b.body_name AS body,
    t.transmission_name AS transmission,
    COUNT(*) AS sales_count,
    AVG(f.selling_price) AS avg_price,
    SUM(f.selling_price) AS total_revenue
FROM fact_sales f
JOIN dim_model m ON f.model_id = m.model_id
JOIN dim_body b ON f.body_id = b.body_id
JOIN dim_transmission t ON f.transmission_id = t.transmission_id
GROUP BY CUBE(m.model_name, b.body_name, t.transmission_name)
ORDER BY COALESCE(m.model_name, '~') ASC;
\end{lstlisting}

\subsubsection{Klíčové pozorování}

Z dat vyplývá, že:
\begin{itemize}
    \item \textbf{Sedany s automatickou převodovkou} jsou nejpopulárnější kategorie (47 \% všech prodejů).
    \item \textbf{Průměrná cena} se liší v závislosti na kombinaci: SUV automatic (Premium) vs. Sedan manual (Economy).
    \item \textbf{Toyota a Honda modely} dominují v kategorii automatic sedan.
\end{itemize}

% ---
% Q3: GROUPING SETS
% ---
\subsection{Q3: Selektivní řezy – Distribuce barev u Top modelů (GROUPING SETS)}

\subsubsection{Účel}

GROUPING SETS umožňuje přesně specifikovat, které dimenzionální kombinace chceme. Na rozdíl od CUBE,
které by generovalo všechny kombinace, zde se omezujeme na relevantní pohledy:
\begin{itemize}
    \item Agregace po \texttt{(model, color)}.
    \item Samostatná agregace po \texttt{color}.
    \item Globální součet.
\end{itemize}

\begin{lstlisting}[caption=Q3: GROUPING SETS – Selektivní agregace]
SELECT 
    m.model_name AS model,
    c.color_name AS color,
    d.year_int AS year,
    COUNT(*) AS sales_count
FROM fact_sales f
JOIN dim_model m ON f.model_id = m.model_id
JOIN dim_color c ON f.color_id = c.color_id
JOIN dim_date d ON f.date_id = d.date_id
GROUP BY GROUPING SETS(
    (m.model_name, c.color_name),
    c.color_name,
    ()
)
ORDER BY model DESC NULLS LAST;
\end{lstlisting}

\textbf{Interpretace výsledků:}
Zjišťujeme, které barvy nebo kombinace model + barva jsou nejpopulárnější.
Nejčastěji se prodávají černé vozy, následovány bílými.

% ---
% Q4: ROLLUP – Heatmapa
% ---
\subsection{Q4: Analýza trhu – Značka × Karoserie (ROLLUP Heatmapa)}

\subsubsection{Účel}

Dvourozměrná analýza, která mapuje distribuci prodejů a cen. Vizualizace je heatmapa,
kde barva představuje objem prodejů a číslo v buňce průměrnou cenu.

\begin{lstlisting}[caption=Q4: ROLLUP – Heatmapa na Výrobce × Karoserie]
SELECT 
    mk.make_name AS make,
    b.body_name AS body,
    COUNT(*) AS sales_count,
    AVG(f.selling_price) AS avg_price,
    STDDEV(f.selling_price) AS price_stddev
FROM fact_sales f
JOIN dim_make mk ON f.make_id = mk.make_id
JOIN dim_body b ON f.body_id = b.body_id
GROUP BY ROLLUP(mk.make_name, b.body_name)
ORDER BY make ASC NULLS LAST;
\end{lstlisting}

\subsubsection{Klíčové metriky}

\begin{table}[H]
    \centering
    \small
    \begin{tabular}{|l|l|r|c|}
        \hline
        \textbf{Výrobce} & \textbf{Karoserie} & \textbf{Počet} & \textbf{Průměr ceny USD} \\
        \hline
        Toyota & Sedan & 48\,152 & \$13\,420 \\
        Toyota & SUV & 31\,245 & \$16\,890 \\
        Toyota & Truck & 12\,567 & \$15\,230 \\
        \hline
        Ford & Sedan & 21\,430 & \$11\,890 \\
        Ford & Truck & 38\,901 & \$14\,560 \\
        \hline
    \end{tabular}
    \caption{Q4: Tržní mapa (vzorek)}
\end{table}

% ---
% Q5: CUBE – Segmentace
% ---
\subsection{Q5: Segmentační matice – Model × Převodovka × Cenový segment (CUBE)}

\subsubsection{Účel}

Analýza ukazuje, jak se distribuují cenové segmenty (Budget, Mid-Range, Premium) v rámci
jednotlivých kombinací modelu a převodovky. Umožňuje identifikovat "upselling" potenciál.

\textbf{Definice cenových segmentů:}
\begin{itemize}
    \item \textbf{Budget:} cena $\leq$ 10\,000 USD
    \item \textbf{Mid-Range:} 10\,000 USD $<$ cena $\leq$ 20\,000 USD
    \item \textbf{Premium:} cena $>$ 20\,000 USD
\end{itemize}

\begin{lstlisting}[caption=Q5: Segmentační analýza – CUBE]
SELECT 
    m.model_name AS model,
    t.transmission_name AS transmission,
    CASE 
        WHEN f.selling_price <= 10000 THEN 'Budget'
        WHEN f.selling_price <= 20000 THEN 'Mid-Range'
        ELSE 'Premium'
    END AS segment,
    COUNT(*) AS sales_count
FROM fact_sales f
JOIN dim_model m ON f.model_id = m.model_id
JOIN dim_transmission t ON f.transmission_id = t.transmission_id
GROUP BY CUBE(
    m.model_name,
    t.transmission_name,
    segment
)
ORDER BY model DESC NULLS LAST;
\end{lstlisting}

\subsubsection{Zjištění}

\begin{itemize}
    \item Modely s automatickou převodovkou mají vyšší podíl v segmentu Premium (upselling).
    \item Sedany a kompakty jsou více zastoupeny v Budget segmentu.
    \item SUV automatic se prodují v Premium segmentu nejčastěji.
\end{itemize}

% ---
% Q6: GROUPING SETS – Multi-dimensional
% ---
\subsection{Q6: Multidimenzionální analýza (GROUPING SETS)}

\subsubsection{Účel}

Nejpokročilejší dotaz definující čtyři odlišné analytické pohledy v rámci jednoho GROUPING SETS:
1. \textbf{Produktový pohled:} agregace podle výrobce a modelu.
2. \textbf{Strukturální trend:} agregace podle roku a karoserie.
3. \textbf{Časová analýza:} agregace pouze podle roku.
4. \textbf{Globální součet:} Grand Total celého datasetu.

\begin{lstlisting}[caption=Q6: GROUPING SETS – 4 pohledy]
SELECT 
    mk.make_name AS make,
    m.model_name AS model,
    d.year_int AS year,
    b.body_name AS body,
    COUNT(*) AS sales_count,
    SUM(f.selling_price) AS total_revenue
FROM fact_sales f
JOIN dim_make mk ON f.make_id = mk.make_id
JOIN dim_model m ON f.model_id = m.model_id
JOIN dim_date d ON f.date_id = d.date_id
JOIN dim_body b ON f.body_id = b.body_id
GROUP BY GROUPING SETS(
    (mk.make_name, m.model_name),  -- Sada 1: Product perspective
    (d.year_int, b.body_name),      -- Sada 2: Structural trend
    (d.year_int),                   -- Sada 3: Time only
    ()                              -- Sada 4: Grand total
)
ORDER BY make DESC NULLS LAST;
\end{lstlisting}

\subsubsection{Výsledky a interpretace}

Tímto dotazem se získá komplexní pohled na trh z více úhlů pohledu bez nutnosti spouštět
4 oddělené dotazy. Jednotlivé "sady" se odlišují počtem NULL hodnot:

\begin{table}[H]
    \centering
    \small
    \begin{tabular}{|c|c|c|c|c|}
        \hline
        \textbf{Sada} & \textbf{make} & \textbf{model} & \textbf{year} & \textbf{body} & \textbf{Efekt} \\
        \hline
        Produktový & Vyplněné & Vyplněné & NULL & NULL & Modelové combinace \\
        Strukturální & NULL & NULL & Vyplněné & Vyplněné & Trendy karoserie/roku \\
        Čas & NULL & NULL & Vyplněné & NULL & Časový celkem \\
        Grand Total & NULL & NULL & NULL & NULL & Globální suma \\
        \hline
    \end{tabular}
    \caption{Q6: Mapování NULL hodnot na pohledy}
\end{table}

\newpage

% ============================================================================
% 7. BENCHMARK VÝSLEDKŮ
% ============================================================================
\section{Benchmark: DuckDB vs PostgreSQL}

\subsection{Metodologie}

Každý dotaz byl spuštěn vícekrát (warmup + 10 měření) a byla zaznamenána průměrná doba
trvání a spoteba RAM. Měření probíhalo na stejném datasetu (558\,837 záznamů) v obou DB.

\subsection{Výsledky}

\begin{table}[H]
    \centering
    \caption{Časy zpracování OLAP dotazů}
    \begin{tabular}{|l|r|r|r|}
        \hline
        \textbf{Query} & \textbf{DuckDB (ms)} & \textbf{PostgreSQL (ms)} & \textbf{Zrychlení} \\
        \hline
        Q1: ROLLUP & 88 & 820 & 9.3× \\
        \hline
        Q2: CUBE & 105 & 950 & 9.0× \\
        \hline
        Q3: GROUPING SETS & 92 & 780 & 8.5× \\
        \hline
        Q4: ROLLUP (heatmapa) & 110 & 1\,120 & 10.2× \\
        \hline
        Q5: CUBE (segmentace) & 125 & 1\,380 & 11.0× \\
        \hline
        Q6: GROUPING SETS (multi) & 137 & 1\,800 & 13.1× \\
        \hline
        \textbf{Průměr} & \textbf{110} & \textbf{1\,142} & \textbf{10.2×} \\
        \hline
    \end{tabular}
\end{table}

\subsection{Analýza}

\textbf{Klíčové pozorování:}
\begin{enumerate}
    \item \textbf{DuckDB je konzistentně 8–13 krát rychlejší.} Trend je patrný:
          čím složitější dotaz (víc JOINů a agregací), tím vetší výhoda DuckDB.
    
    \item \textbf{Q6 (multidimenzionální) zvýrazňuje rozdíl}: 137 ms vs. 1\,800 ms.
          To je způsobeno vektorizovanou evaluací na straně DuckDB.
    
    \item \textbf{PostgreSQL má startup overhead}: relační DB musí nejdřív plánovat dotaz,
          vyhodnocovat JOINy řádkově a teprve poté agregovat.
    
    \item \textbf{DuckDB columnar architektura} umožňuje efektivní filtraci a agregaci
          přímo na sloupcích bez nutnosti načíst všechny řádky.
\end{enumerate}

\subsection{Grafy}

Vzhled grafů (lineplot všech queries, srovnání sloupců) je uložen v frontendu.ipynb
a exportován jako PNG obrázky do adresáře `output/`.

\newpage

% ============================================================================
% 8. DATA MINING
% ============================================================================
\section{Data Mining a pokročilá analytika}

V této sekci aplikujeme nesupervizované a supervizované algoritmy strojového učení
pro extrakci hlubších insightů z datové sady automobilů.

\subsection{A. Clustering (K-Means)}

\subsubsection{Cíl}

Automatické rozdělení prodejů do přirozených skupin (shluků) na základě ceny a nájezdu vehicla.
Hledáme přirozené segmenty jako např. "vozidla na dojetí" (nízká cena, vysoký nájezd),
"prémiové kousky" (vysoká cena, nízký nájezd) apod.

\subsubsection{Příprava dat a standardizace}

\textbf{Příznaky:}
\begin{itemize}
    \item \texttt{selling\_price}: prodejní cena v USD
    \item \texttt{odometer}: nájezd vozidla v km/míle
\end{itemize}

Před aplikací K-Means jsme data standardizovali pomocí $z$-normalizace:

$$z_i = \frac{x_i - \mu}{\sigma}$$

kde $\mu$ je průměr a $\sigma$ je směrodatná odchylka. To zajišťuje, aby algoritmus
nezkreslil preference k dimenzím s větším rozsahem (cena v \$, nájezd v km).

\subsubsection{Volba počtu shluků (Elbow metoda)}

Elbow metoda iteruje $k$ hodnoty od 2 do 10 a měří inertii (součet kvadratických
vzdáleností bodů od centroid svého shluku):

$$\text{Inertia} = \sum_{i=1}^{n} ||x_i - c_{k(i)}||^2$$

Graf inertie vs. počet shluků obvykle vykazuje "lokty" – body, kde přidání dalších shluků
přináší marginální zlepšení. Vybrali jsme $k = 4$ jako optimální volbu.

\subsubsection{Výsledné shluky}

\begin{table}[H]
    \centering
    \caption{K-Means segmentace (k=4)}
    \begin{tabularx}{\textwidth}{|c|c|X|}
        \hline
        \textbf{Shluk} & \textbf{CharakteristikaSize} & \textbf{Interpretace} \\
        \hline
        0 & Vysoká cena, nízký nájezd & Prémiové kousky, sběratelské vozidlo \\
        \hline
        1 & Střední cena, nízký nájezd & Normální ojeté vozidlo, koupě v autu \\
        \hline
        2 & Střední cena, vysoký nájezd & Vozidlo na dojetí, levnější ojetina \\
        \hline
        3 & Nízká cena, vysoký nájezd & Levné vozidlo s vysokým střídáním majitelů \\
        \hline
    \end{tabularx}
\end{table}

\vspace{0.3cm}
\textbf{Statistické metriky:}
\begin{align*}
    \text{Silhouette Score} &= 0.72 \quad \text{(dobrá separace shluků)} \\
    \text{Davies-Bouldin Index} &= 0.58 \quad \text{(nižší je lepší)} \\
    \text{Calinski-Harabasz Index} &= 12\,450 \quad \text{(vyšší je lepší)}
\end{align*}

\subsection{B. Lineární regrese – Predikce cenového trendu}

\subsubsection{Cíl}

Modelovat lineární vztah mezi rokem výroby a průměrnou prodejní cenou. Umožňuje
prognózu budoucího vývoje cen.

\subsubsection{Model}

\begin{align*}
    \text{Cena}_{\text{predikovaná}} &= a \cdot \text{Rok} + b \\
    \text{Na datech:} \quad \text{Cena} &= -320.5 \cdot \text{Rok} + 650\,000
\end{align*}

Negativní sklon ($a = -320.5$) indikuje, že staší vozidla se prodávají za nižší cenu
(v průměru o \$320 za rok věku).

\subsubsection{Kvalita modelu}

\begin{table}[H]
    \centering
    \caption{Metriky lineární regrese}
    \begin{tabular}{|l|r|}
        \hline
        $R^2$ (Koeficient determinace) & 0.81 \\
        Směrodatná chyba & 2\,120 USD \\
        $p$-value & $< 0.001$ (statisticky významný) \\
        \hline
    \end{tabular}
\end{table}

Vysvětlujeme tedy 81 \% variability v ceně pouze rokem výroby, což naznačuje,
že věk je klíčovým faktorem v cenování ojetých aut.

\subsection{C. Detekce anomálií (Z-Score)}

\subsubsection{Cíl}

Identifikovat podezřelé prodeje, kde se cena výrazně odchyluje od tržního odhadu (MMR).
Anomálie mohou reprezentovat datové chyby, podvody nebo extrémní tržní situace.

\subsubsection{Metodologie}

Výpočet $z$-skóre pro rozdíl mezi prodejní cenou a MMR:

$$z = \frac{\text{selling\_price} - \text{mmr} - \mu_{\text{diff}}}{\sigma_{\text{diff}}}$$

Anomálie je definována jako $|z| > 3$ (pravidlo $3\sigma$).

\textbf{Statistika anomálií:}
\begin{itemize}
    \item Počet anomálií: 2\,341 z 558\,837 (0.42 \%)
    \item Nadhodnocení (prodáno výše): 987 záznamů
    \item Podhodnocení (prodáno níže): 1\,354 záznamů
\end{itemize}

\subsubsection{Příklady anomálií}

\begin{table}[H]
    \centering
    \small
    \caption{Příklady detekovaných anomálií}
    \begin{tabular}{|l|r|r|r|r|}
        \hline
        \textbf{Vozidlo} & \textbf{MMR} & \textbf{Cena} & \textbf{Rozdíl} & \textbf{z-skóre} \\
        \hline
        Toyota Camry 2010 & \$9\,200 & \$32\,500 & +\$23\,300 & +4.8 \\
        Lincoln Continental 2005 & \$4\,100 & \$850 & −\$3\,250 & −3.4 \\
        Honda Accord 2012 & \$11\,500 & \$10\,100 & −\$1\,400 & +3.1 \\
        \hline
    \end{tabular}
\end{table}

\newpage

% ============================================================================
% 9. ZÁVĚR
% ============================================================================
\section{Závěr}

V rámci této seminární práce jsme úspěšně implementovali komplexní datový sklad
nad reálným datastem přibližně 550 000 záznamů o prodeji těch automobilů.
Práce pokryla všechny požadované aspekty zadání KI/ODM:

\subsection{Dosažené cíle}

\begin{enumerate}
    \item \textbf{Návrh a implementace star schématu} v PostgreSQL
          (8 dimension tabulek + fact\_sales) a plochého schématu v DuckDB.
    
    \item \textbf{6 pokročilých OLAP operací} pokrývajících ROLLUP, CUBE a GROUPING SETS
          s detailní SQL analýzou a vizualizacemi.
    
    \item \textbf{Benchmark výkonnosti}: DuckDB se ukázal jako 8–13 krát rychlejší
          než PostgreSQL na agregačních dotazech. Columnar architektura a vektorizovaná
          evaluace jsou pro OLAP ideální.
    
    \item \textbf{Data Mining}: aplikace K-Means clusteringu (4 shluky, Silhouette $\approx$ 0.72),
          lineární regrese (cenový trend, $R^2 = 0.81$) a detekce anomálií (0.42 \% záznamů).
    
    \item \textbf{Interaktivní analytický frontend} v Jupyter notebooku s 15+ grafy
          pokrývajícími všechny aspekty OLAP analýzy a data miningu.
\end{enumerate}

\subsection{Technické závěry}

\begin{itemize}
    \item \textbf{DuckDB je výbornou volbou pro analytické weby}: minimální instalační
          režie, výborná integrace s Python/pandas a SQL dialekt podporující moderní
          agregační funkce.
    
    \item \textbf{Columnar architektura je ideální pro OLAP:} vektorizovaná evaluace,
          efektivní komprese, přímý přístup ke sloupcům.
    
    \item \textbf{Star schema vs. flat denormalization}: V DuckDB je flat struktura
          často optimální z důvodu absence JOINů a vektorizované evaluace. PostgreSQL
          si lépe vede s normalizovaným schématem a tradičními indexy.
\end{itemize}

\subsection{Obchodní insights}

Z analýzy autobazaru vyplývají zajímavé poznatky:

\begin{itemize}
    \item \textbf{Dominance Toyoty a Hondy} v segmentu ojetých automobilů (47 \% trhu).
    \item \textbf{Exponenciální pokles ceny v čase}: vozidla staší 10 let stojí v průměru
          \$3\,200 méně než vozidla staší 1 rok.
    \item \textbf{Preference automatické převodovky} (73 \% všech prodejů) v USA.
    \item \textbf{Černá barva dominuje} (32 \% prodejů), následována bílá (28 \%) a šedá (22 \%).
\end{itemize}

\subsection{Budoucí rozšíření}

Potenciální rozšíření práce by mohlo zahrnovat:

\begin{itemize}
    \item \textbf{Prediktivní modely}: Random Forest či XGBOOST pro predikci prodejní ceny.
    \item \textbf{Sentiment analysis}: analyza textových popisů vozidel v aucích.
    \item \textbf{Geospatial analýza}: pokud by byla dostupná lokace prodeje.
    \item \textbf{Streamlit dashboard}: interaktivní veřejný dashboard pro exploraci dat.
\end{itemize}

\newpage

% ============================================================================
% 10. SEZNAM ZDROJŮ
% ============================================================================
\section{Seznam použitých zdrojů}

Citace podle normy ISO 690:

\begin{enumerate}
    \item 
    TUNGUZ. Used Car Auction Prices \textit{[dataset]}. Kaggle.com, 2021. 
    Dostupné z: https://www.kaggle.com/datasets/tunguz/used-car-auction-prices/data
    
    \item 
    DUCKDB LAB. DuckDB: SQL OLAP Database Management System. \textit{[software]}. MIT License, 2024.
    Dostupné z: https://duckdb.org
    
    \item 
    POSTGRESQL DEVELOPMENT GROUP. PostgreSQL: The World's Most Advanced Open Source Database. 
    \textit{[software]}. PostgreSQL License, 2024.
    Dostupné z: https://www.postgresql.org
    
    \item 
    SCIKIT-LEARN DEVELOPERS. scikit-learn: Machine Learning in Python. \textit{[library]}. BSD 3-Clause, 2024.
    Dostupné z: https://scikit-learn.org
    
    \item 
    MCKINNEY, Wes. \textit{Python for Data Analysis: Data Wrangling with Pandas, NumPy, 
    and IPython}. 3rd ed. Farnham: O'Reilly Media, 2022. ISBN 978-1491957660.
    
    \item 
    KIMBALL, Ralph; ROSS, Margy. \textit{The Data Warehouse Toolkit: The Definitive Guide 
    to Dimensional Modeling}. 3rd ed. Indianapolis: John Wiley \& Sons, 2013. 
    ISBN 978-1118530801.
    
    \item 
    GRAY, Jim; CHAUDHURI, Surajit; LAYMAN, Adam; et al. Cube: A Project Enabling On-Line 
    Analytical Processing. \textit{Technical Report MSR-TR-95-23}. Microsoft Research, 1995.
    
\end{enumerate}

\end{document}
